% Part: first-order-logic
% Chapter: lindstrom
% Section: abstract-logics

\documentclass[../../../include/open-logic-section]{subfiles}

\begin{document}

\olfileid{mod}{lin}{alg}
\olsection{تجریدی منطقاں}

\begin{defn}
اک \emph{تجریدی منطق} اک جوڑا $\tuple{L, \models_L}$ اے؛ اِتھے $L$
اوہ فنکشن اے جو ہر !!{language}~$\Lang{L}$ نوں
!!{sentence}s دا سیٹ $L(\Lang{L})$ دیندا اے، تے $\models_L$
اُس !!{language}~$\Lang{L}$ دیاں !!{structure}s تے
$L(\Lang{L})$ دے !!{element}s وچکار اک رشتہ اے۔ خاص طور اُتے
$\tuple{F, \models}$ عام پہلے درجے دی منطق اے: $F$ اوہ فنکشن اے
جو !!{language}~$\Lang{L}$ نوں $\Lang{L}$ دیاں غیر منطقی علامتاں توں بنے
پہلے درجے دے !!{sentence}s دا سیٹ دیندا اے، تے $\models$
پہلے درجے دی منطق دا صدقی رشتہ اے۔
\end{defn}

دھیان دیو کہ کسے مقرر !!{language} لئی اسیں ہن وی
!!{structure} دا اوہو مفہوم ورتدے آں جو پہلے درجے دی منطق وچ
سی۔ البتہ اسیں ایہ پہلے توں نہیں مندے کہ !!{sentence}s عام
طریقے نال $\Lang{L}$ دیاں بنیادی علامتاں توں بنے ہون، نہ ایہ
کہ رشتہ $\models_L$ پہلے درجے دی منطق وانگ تکراری تعریف نال
مقرر ہوندا ہووے۔ مثلاً ایہ تعریف پوری طرح عمومی اے؛ ایس دا
مقصد $\tuple{L,\models_L}$ وچ ایسے !!{sentence}s نوں وی شامل
کرنا اے جنہاں وچ لامتناہی لمبے
عطف یا فصل، یا $\lexists$ تے $\lforall$ توں ہٹ کے مسوّر
(مثلاً «لامتناہی طور تے بہت سارے~$x$ ایس طرح دے نیں کہ \dots»)،
یا شاید مسوّراں دے لامتناہی لمبے ابتدائی سلسلے ہون۔ ایہ گل
صاف رکھن لئی کہ $L(\Lang{L})$ دے «!!{sentence}s» لازماً عام
پہلے درجے دے !!{sentence}s نہیں، ایس باب وچ اسیں اوہناں لئی
!!{variable}s $!E$, $!F$,~\dots ورتاں گے، تے $!A$, $!B$,~\dots
صرف عام پہلے درجے دے !!{formula}s لئی رکھاں گے۔

\begin{defn}
$\Mod(L){!E}$ توں مراد طبقہ $\Setabs{\Struct{M}}{\Struct{M}
  \models_L !E}$ لوو۔ جے !!{language} صراحت نال دسنی ہووے،
تاں اسیں $\Mod[L](L){!E}$ لکھدے آں۔ دو
!!{structure}s $\Struct{M}$ تے $\Struct{N}$،
$\Lang{L}$ لئی،
$\tuple{L, \models_L}$ وچ \emph{جملہ وار ہم ارز} ہوندیاں نیں،
جِنہاں نوں $\Struct{M} \elemequiv[L] \Struct{N}$ لکھدے آں،
جے $L(\Lang{L})$ دے اوہو !!{sentence}s دوہاں وچ سچ ہون۔
\end{defn}

\begin{defn}
اک تجریدی منطق $\tuple{L,\models_L}$،
!!{language} $\Lang{L}$ لئی \emph{معمول مطابق (normal)} اے جے ایہ
ہیٹھاں لکھیاں خاصیتاں پوری کرے:
\begin{enumerate}
\item (\emph{$L$-یک رُخیت}) !!{language}s $\Lang{L}$ تے
  $\Lang{L'}$ لئی، جے $\Lang{L} \subseteq \Lang{L'}$، تاں
  $L(\Lang{L}) \subseteq L(\Lang{L'})$۔
\item (\emph{توسیع دی خاصیت}) ہر $!E \in L(\Lang{L})$ لئی
  اک \emph{محدود} ذیلی سیٹ $\Lang{L'}$ موجود اے، جو
  $\Lang{L}$ دا اے،
  ایس طرح کہ رشتہ $\Struct{M} \models_L !E$ صرف
  $\Struct{M}$ دے $\Lang{L'}$ تک گھٹائے روپ اُتے منحصر ہووے؛
  یعنی جے $\Struct{M}$ تے $\Struct{N}$ دا $\Lang{L'}$ تک
  گھٹایا روپ اکّو ہووے، تاں $\Struct{M}
  \models_L !E$ جے تے صرف جے $\Struct{N} \models_L !E$۔
\item (\emph{ہم شکلی دی خاصیت}) جے $\Struct{M} \models_L !E$
  تے $\Struct{M} \simeq \Struct{N}$، تاں
  $\Struct{N} \models_L !E$ وی۔
\item (\emph{نواں نام رکھن دی خاصیت}) رشتہ $\models_L$ نام
  بدلنے نال محفوظ رہندا اے: جے !!{language} $\Lang{L}'$
  $\Lang{L}$ دے ہر نشان $P$ دی تھاں اوہو تعدادِ دلائل
  والا نشان $P'$ تے ہر مستقل $c$ دی تھاں وکھرا مستقل $c'$
  رکھ کے بنے، تاں ہر !!{structure}~$\Struct{M}$ تے
  !!{sentence}~$!E$ لئی $\Struct{M}
  \models_L !E$ جے تے صرف جے $\Struct{M}' \models_L !E'$،
  جتھے $\Struct{M}'$ اک $\Lang{L}'$-!!{structure} اے جو
  $\Lang{L}$ دی اصل ساخت دے اکو دائرے تے تعبیراں نال بن دی اے؛
  صرف علامتاں دے ناں بدل دے نیں، تے سارے نویں ناں آپس وچ وکھرے نیں۔
  ایتھے $!E' \in L(\Lang{L}')$۔
\item (\emph{بولین خاصیت}) تجریدی منطق $\tuple{L,
  \models_L}$ بولین رابطیاں دے تحت ایس معنی وچ بند اے
  کہ ہر $!E \in L(\Lang{L})$ لئی کوئی~$!F \in
  L(\Lang{L})$ موجود اے جس لئی $\Struct{M} \models_L !F$
  جے تے صرف جے $\Struct{M} \not\models_L !E$، تے ہر
  $!E$ تے $!F$ لئی کوئی $!G$ موجود اے جس لئی
  $\Mod(L){!G} = \Mod(L){!E} \cap
  \Mod(L){!F}$۔ ایس طرح بند جوہری !!{formula}s تے ہور
  رابطیاں لئی وی۔
\item (\emph{مسوّر دی خاصیت}) ہر مستقل $c$ جو $\Lang{L}$ وچ اے
  تے $!E \in L(\Lang{L})$ لئی کوئی $!F \in L(\Lang{L'})$
  موجود اے ایس طرح کہ
  \[
  \Mod[L'](L){!F} = \Setabs{\Struct{M}}{\Expan{M}{a} \in
  \Mod[L](L){!E} \text{ کسے } a \in \Domain{M}\text{ لئی}},
  \]
  جتھے $\Lang{L'} = \Lang{L} \setminus \{c\}$، تے
  $\Expan{M}{a}$، $\Struct{M}$ دی $\Lang{L}$ تک اوہ توسیع
  اے جس وچ $a$، $c$ نوں دتی قدر اے۔
\item (\emph{نسبتی بنانے دی خاصیت}) جے !!a{sentence} $!E \in
  L(\Lang{L})$ تے نشان $R$, $c_1$, \dots, $c_n$،
  $\Lang{L}$ وچ نہ ہون، تاں کوئی !!a{sentence} $!F \in
  L(\Lang{L} \cup \{R,c_1,\ldots,c_n\})$ موجود اے، جس نوں
  $!E$ دا $\Atom{R}{x, c_1, \dots c_n}$ تک
  \emph{نسبتی روپ} آکھدے نیں، ایس طرح کہ ہر
  !!{structure}~$\Struct{M}$ لئی، جدوں منتخب دائرہ غیر خالی ہووے
  تے اصل زبان دے ہر مستقل دی قدر ایس وچ ہووے۔ ایس باب دیاں
  تعلقی زباناں لئی ایہ شرطاں اوہ دائرہ ذیلی ساخت بناون لئی لازم نیں:
  \[
  \Expan{M}{X, b_1, \ldots, b_n} \models_L !F \text{ اوہدوں تے
    صرف اوہدوں جدوں } \Struct{N} \models_L !E,
  \]
  جتھے $\Struct{N}$، $\Struct{M}$ دی اوہ ذیلی ساخت اے
  جس دا !!{domain}
  $\Domain{N} = \Setabs{a\in \Domain{M}}{\tuple{a, b_1, \dots, b_n} \in X}$
  اے (ویکھو \olref[bas][sub]{rem:substructure})، تے
  $\Expan{M}{X, b_1, \ldots,
    b_n}$، $\Struct{M}$ دی اوہ توسیع اے جو بالترتیب
  $R$, $c_1$, \dots, $c_n$ دی تعبیر
  $X$, $b_1,$ \dots, $b_n$ نال کردی اے (جتھے
  $X \subseteq M^{n+1}$)۔
\end{enumerate}
\end{defn}
 
\begin{defn}
دو تجریدی منطقاں $\tuple{L_1, \models_{L_1}}$ تے
$\tuple{L_2, \models_{L_2}}$ دتیاں ہون۔ اسیں آکھدے آں کہ
دوجی منطق پہلی دے \emph{گھٹ توں گھٹ برابر اظہار جوگی} اے،
جس نوں $\tuple{L_1, \models_{L_1}}
\leq \tuple{L_2, \models_{L_2}}$ لکھدے آں، جے ہر
!!{language} $\Lang{L}$ تے ہر !!{sentence}
$!E \in L_1(\Lang{L})$ لئی کوئی !!a{sentence} $!F
\in L_2(\Lang{L})$ موجود ہووے جس لئی
$\Mod[L](L_1){!E} =
\Mod[L](L_2){!F}$۔ منطقاں
$\tuple{L_1, \models_{L_1}}$ تے
$\tuple{L_2, \models_{L_2}}$ \emph{ہم ارز} نیں جے
$\tuple{L_1,
  \models_{L_1}} \leq \tuple{L_2, \models_{L_2}}$ تے
$\tuple{L_2,
  \models_{L_2}} \leq \tuple{L_1, \models_{L_1}}$۔
\end{defn}

\begin{rem}
  پہلے درجے دی منطق، یعنی تجریدی منطق $\tuple{F, \models}$،
  معمول مطابق اے۔ دراصل اُتّے دتیاں خاصیتاں پہلے درجے دی
  منطق لئی زیادہ تر سیدھیاں نیں۔ صرف ایہ یاد رکھو کہ توسیع
  دی خاصیت غیر متعلقہ علامتاں توں آزادی
  (Extensionality/Relevance) اُتے آ ٹکدی اے، تے !!a{sentence} $!E$
  دا $\Atom{R}{x, c_1, \dots, c_n}$ تک نسبتی روپ لبھن لئی
  پہلاں وجودی مسوّراں نوں نفی تے عمومی مسوّر دے مساوی روپ وچ لکھو؛ فیر
  ہر !!{subformula} $\lforall[x][!F]$ دی تھاں
  $\lforall[x][(\Atom{R}{x, c_1, \dots, c_n} \lif \ !F)]$
  رکھ کے لبھدا اے۔ ہور ایہ کہ جے
  $\tuple{L, \models_L}$ معمول مطابق ہووے، تاں
  $\tuple{F, \models} \leq \tuple{L, \models_{L}}$،
  جیویں پہلے درجے دے !!{formula}s اُتے استقرا نال
  دکھایا جا سکدا اے۔ ایس لئی عمومیت گنوائے بغیر اسیں
  فرض کر سکدے آں کہ ہر پہلے درجے دا !!{sentence} ہر
  معمول مطابق منطق وچ شامل اے۔
\end{rem}

\end{document}
